\documentclass[11pt,a4paper]{article}

\usepackage[margin=2.5cm]{geometry}
\usepackage{amsmath,amssymb,amsthm}
\usepackage{graphicx}
\usepackage{booktabs}
\usepackage{array}
\usepackage{hyperref}
\usepackage{microtype}

\newtheorem{theorem}{Theorem}[section]

\newcommand{\R}{\mathbb{R}}
\newcommand{\dd}{\mathrm{d}}
\newcommand{\topT}{^{\top}}
\newcommand{\grad}{\nabla}
\newcommand{\Div}{\nabla\!\cdot\!}
\newcommand{\curl}{\nabla\!\times\!}
\newcommand{\perpgrad}{\nabla^{\!\perp}}
\newcommand{\Drift}{D^{(1)}}
\newcommand{\Diff}{D^{(2)}}
\newcommand{\norm}[1]{\lVert #1 \rVert}

\title{Reconstruction of Scalar Potentials and Rotational Components\\
from Latent Spaces with P-spline Kramers--Moyal Estimation}
\author{}
\date{November 11, 2026}

\begin{document}

\maketitle

\begin{abstract}
We present a fully data-driven pipeline for reconstructing effective scalar
potentials and rotational components from low-dimensional latent trajectories
extracted from EEG recordings. Drift and diffusion fields are estimated from
the Kramers--Moyal conditional moments by \emph{global penalised B-spline
regression} (P-splines), which replaces the local Nadaraya--Watson kernel
estimator and removes the need for spatial binning at the estimation stage.
The diffusion matrix is inverted through adaptive SVD truncation, and both the
conservative and solenoidal parts of the effective field are recovered using a
variational Galerkin--B-spline formulation (Isogeometric Analysis) with
built-in Helmholtz decomposition. In two dimensions the decomposition decouples
into two independent scalar Poisson equations---one for the potential and one
for the stream function---solved with identical stiffness matrices. A
non-equilibrium metric quantifies the relative weight of the rotational
component. The penalisation parameter $\lambda$ of the P-spline estimator plays
the role of the kernel bandwidth of the classical estimator and becomes the
central hyperparameter to calibrate; we define the full hyperparameter set of
the method and a validation strategy to select it. Numerical experiments on
equilibrium and non-equilibrium stochastic systems are reported in
Section~\ref{sec:experiments} [placeholders pending final runs].
\end{abstract}

\tableofcontents
\newpage

%==============================================================================
\section{Introduction}
%==============================================================================

We work with time-series data of dimensionality 1 to 5 that represent a latent
space extracted from EEG recordings. The extraction pipeline combines PCA
whitening, Independent Component Analysis (ICA), Markov-time estimation, and
the computation of conservative fractions to obtain a low-dimensional embedding
that captures the essential stochastic dynamics of the underlying neural
activity \cite{tabar2019}. The question we address is whether, from these
latent trajectories, one can reconstruct an effective scalar potential whose
geometry reveals the attractors, barriers, and basins that govern the system's
behaviour, and, when the system operates far from detailed balance, whether one
can additionally recover the rotational component that drives circulating
currents and limit-cycle dynamics.

The goal is to develop a fully data-driven reconstruction pipeline that takes
the latent time series, estimates the drift and diffusion fields from the
Kramers--Moyal conditional moments \cite{risken1989}, and recovers both the
conservative and solenoidal parts of the effective field through a variational
Galerkin--B-spline formulation (Isogeometric Analysis, see
\cite{cottrell2009}) with built-in Helmholtz decomposition. The resulting
potentials and stream functions can then be catalogued and compared across
subjects or brain states, opening the door to a ``zoology of potentials'' that
characterises different dynamical regimes in an interpretable, low-dimensional
form.

The present document extends the pipeline described in the companion report
(\emph{Reconstruction of Scalar Potentials and Rotational Components from
Latent Spaces}) by replacing the estimation stage. In the original pipeline the
Kramers--Moyal coefficients are estimated locally by Nadaraya--Watson kernel
regression on a regular grid of bins, which requires choosing a kernel, a
bandwidth, a bin count, a post-smoothing scale, and a density threshold, and
which suffers from boundary bias and from the curse of dimensionality through
cell occupancy. Here the coefficients are instead estimated by
\emph{P-splines} \cite{eilers1996}: a global regression on a tensor-product
B-spline basis with a difference penalty on the coefficients. The estimator is
defined without any spatial discretisation of the data, is naturally robust to
boundary effects through the penalty, and reduces the bandwidth selection
problem to the calibration of a scalar penalisation parameter $\lambda$ (one
for the drift, one for the diffusion). All downstream stages---robust
inversion of the diffusion matrix, IgA reconstruction, Helmholtz decomposition,
and the non-equilibrium metric---are unchanged, since the P-spline estimator
delivers the same outputs (drift field, diffusion field, grid edges) evaluated
on an arbitrary output grid.

A key motivation for extending the method beyond the scalar potential is that
many neural and physical systems exhibit non-equilibrium dynamics---circulating
probability currents, limit cycles, and oscillatory behaviour---that a purely
conservative potential cannot reproduce. In such systems the drift field
contains a solenoidal component that carries essential dynamical information.
By reconstructing both components, the pipeline provides a complete picture of
the effective dynamics directly from data.

%==============================================================================
\section{Physical framework and problem formulation}
%==============================================================================

\subsection{Fokker--Planck equation and Kramers--Moyal coefficients}

Consider a $D$-dimensional diffusion process governed by the It\^o stochastic
differential equation
\begin{equation}
\dd X_t = f(X_t)\,\dd t + \sigma(X_t)\,\dd W_t,
\label{eq:sde}
\end{equation}
where $f:\R^D \to \R^D$ is the drift, $\sigma:\R^D \to \R^{D\times M}$ is the
noise coupling matrix, and $W_t$ is an $M$-dimensional Wiener process. The
probability density $\rho(x,t)$ evolves according to the forward
Fokker--Planck equation \cite{risken1989,gardiner2009}
\begin{equation}
\frac{\partial \rho(x,t)}{\partial t}
= -\Div \big[ f(x)\,\rho(x,t) \big]
+ \grad\grad : \big[ D(x)\,\rho(x,t) \big],
\label{eq:fpe}
\end{equation}
where $D(x) = \tfrac{1}{2}\sigma(x)\sigma(x)\topT$ is the symmetric positive
semi-definite diffusion matrix. The coefficients that appear in the
Kramers--Moyal expansion are the conditional moments
\begin{equation}
\Drift_i(x) = \lim_{\Delta t \to 0}
\frac{\mathbb{E}[\Delta X_i \mid X_t = x]}{\Delta t},
\qquad
\Diff_{ij}(x) = \lim_{\Delta t \to 0}
\frac{\mathbb{E}[\Delta X_i\, \Delta X_j \mid X_t = x]}{2\,\Delta t},
\label{eq:km_moments}
\end{equation}
which coincide with $f_i(x)$ and $D_{ij}(x)$, respectively. Pawula's theorem
guarantees that the expansion may be truncated at second order whenever $D(x)$
is positive definite \cite{risken1989}.

\subsection{Scalar potential and the non-equilibrium effective field}

In steady state the probability current
\begin{equation}
J(x) = f(x)\,\rho(x) - \Div \big[ D(x)\,\rho(x) \big]
\label{eq:current}
\end{equation}
is solenoidal ($\Div J = 0$). Under the hypothesis of detailed balance one has
$J \equiv 0$, which implies a purely conservative decomposition $f = -D\,\grad
U$ with $U$ a scalar potential; the stationary density then takes the
Boltzmann-like form $\rho \propto \exp(-U)$. Experimental data---in particular
latent spaces extracted from EEG---rarely satisfy detailed balance exactly.

In the general non-equilibrium case the drift decomposes as
\begin{equation}
f(x) = -D(x)\,\grad U(x) + v(x),
\label{eq:drift_decomp}
\end{equation}
where $U$ is the scalar potential and $v$ is a non-conservative force that
encodes circulating currents, limit-cycle dynamics, and other genuinely
non-equilibrium phenomena.

We define the effective field
\begin{equation}
g(x) = -D(x)^{-1} f(x) = \grad U(x) - D(x)^{-1} v(x),
\label{eq:effective_field}
\end{equation}
which reduces to $\grad U$ when $v \equiv 0$. Crucially, $g$ is obtained by a
purely algebraic transformation of the estimated drift and diffusion; no
assumption of detailed balance is made at any stage.

\subsection{Helmholtz decomposition and variational formulation}
\label{sec:helmholtz}

Given the estimated field $g$ on a domain $\Omega$ with sufficient empirical
density, we decompose it into its conservative and solenoidal constituents. The
space of square-integrable vector fields admits the orthogonal Helmholtz
decomposition \cite{strang1986,helmholtz1858}
\begin{equation}
L^2(\Omega;\R^D) = \mathcal{G} \oplus \mathcal{S},
\end{equation}
with
\begin{equation}
\mathcal{G} = \{ \grad u : u \in H^1(\Omega) \},
\qquad
\mathcal{S} = \{ r \in L^2(\Omega;\R^D) : \Div r = 0,\;
r \cdot n|_{\partial\Omega} = 0 \}.
\end{equation}

The conservative component $\grad U$ is obtained by minimising the $L^2$
discrepancy \cite{strang1986}:
\begin{equation}
U^* = \arg\min_{U \in H^1(\Omega)}
\int_\Omega \norm{\grad U(x) - g(x)}^2 \, \dd V .
\label{eq:variational}
\end{equation}
The optimality condition $\delta J = 0$ for arbitrary variations
$\delta U \in H^1(\Omega)$ yields, after integration by parts, the field
equation and natural Neumann boundary condition
\begin{equation}
-\grad^2 U = -\Div g \quad \text{in } \Omega,
\qquad
\frac{\partial U}{\partial n} = g \cdot n \quad \text{on } \partial\Omega .
\label{eq:poisson_U}
\end{equation}
Because $U$ is defined only up to an additive constant, we anchor it at the
point of maximum empirical density,
\begin{equation}
U(x_{\max}) = 0,
\qquad
x_{\max} = \arg\max_{x \in \Omega} \rho(x),
\label{eq:anchoring}
\end{equation}
which corresponds to the theoretical minimum of the potential for stationary
dynamics.

\subsection{Non-equilibrium metric}

The solenoidal residual $r = g - \grad U$ is computed explicitly after solving
for $U$. We define the non-equilibrium metric
\begin{equation}
\eta = \frac{\norm{r}_{L^2,w}}{\norm{g}_{L^2,w}}
= \sqrt{\frac{\sum_k w_k\, \norm{r_k}^2}{\sum_k w_k\, \norm{g_k}^2}},
\label{eq:eta}
\end{equation}
where the norm is weighted by the cell volume and the empirical density. The
physical interpretation is direct:
\begin{itemize}
\item If $\eta < 0.1$, the system is close to detailed balance and $U$ is a
genuine physical potential, interpretable in terms of energy barriers and
Kramers escape rates \cite{risken1989}.
\item If $0.1 \le \eta < 0.5$, the rotational component is moderate: $U$
captures the restoring dynamics while $r$ contains the inertial or circulating
motion.
\item If $\eta \ge 0.5$, the reconstruction is a mathematically valid
projection that serves as a topological descriptor of the stationary
density---to first approximation $U \approx -\log\rho$---but should not be
interpreted as mechanical energy of the system nor used to reconstruct the
full drift via $f = -D\,\grad U$.
\end{itemize}

%==============================================================================
\section{Recovery of the rotational component}
%==============================================================================

\subsection{Stream function and decoupled Poisson system in 2D}

In two dimensions the Helmholtz decomposition admits a particularly convenient
representation. A classical result due to Auchmuty \cite{auchmuty2016}
guarantees that every solenoidal field $r \in \mathcal{S}$ on a simply
connected domain $\Omega \subset \R^2$ with $r \cdot n|_{\partial\Omega} = 0$
can be written as the perpendicular-gradient of a scalar stream function
$\psi$:
\begin{equation}
r(x) = \perpgrad \psi(x) \triangleq
\begin{pmatrix} -\partial_{x_2}\psi \\ \partial_{x_1}\psi \end{pmatrix}.
\label{eq:stream_function}
\end{equation}
The full decomposition of the effective field therefore reads
\begin{equation}
g(x) = \grad U(x) + \perpgrad \psi(x).
\label{eq:hodge_2d}
\end{equation}

Applying the divergence operator to \eqref{eq:hodge_2d} and noting that
$\Div(\perpgrad\psi) = 0$ recovers the Poisson equation for $U$ derived in
Section~\ref{sec:helmholtz}. Applying the scalar curl operator
$\curl g = \partial_{x_1} g_2 - \partial_{x_2} g_1$ and using
$\curl(\grad U) = 0$ yields a second, independent Poisson equation for the
stream function:
\begin{equation}
-\grad^2 \psi = \curl g \quad \text{in } \Omega,
\qquad
\frac{\partial \psi}{\partial n} = g \cdot \tau \quad \text{on } \partial\Omega,
\label{eq:poisson_psi}
\end{equation}
where $\tau = (-n_2, n_1)$ is the unit tangent vector. Both boundary
conditions are natural Neumann conditions that emerge from the variational
formulation. The critical advantage of this 2D formulation is that the system
fully decouples: $U$ and $\psi$ satisfy two independent scalar Poisson
equations, each amenable to the same Galerkin--B-spline machinery.

\subsection{Variational formulation for the stream function}

The weak form for $\psi$ mirrors that of $U$. We seek $\psi \in H^1(\Omega)$
such that
\begin{equation}
\int_\Omega \grad\psi \cdot \grad\chi \, \dd V
= \int_\Omega g \cdot \perpgrad\chi \, \dd V,
\qquad \forall \chi \in H^1(\Omega),
\label{eq:weak_psi}
\end{equation}
where $\perpgrad\chi = (-\partial_{x_2}\chi, \partial_{x_1}\chi)$. This form
does not require numerical differentiation of $g$---the perpendicular-gradient
acts on the test function, not on the data---and preserves the same symmetric
structure as the weak form for $U$. The Galerkin discretisation with B-splines
therefore produces a linear system $K\,\mathbf{d} = F_\psi$ where the stiffness
matrix $K$ is identical to that of the $U$ system and only the load vector
differs: if $F_{U,i} = \sum_{d=1}^{D} \int_\Omega w(x)\, g_d(x)\,
\partial N_i/\partial x_d \, \dd V$, then
\begin{equation}
F_{\psi,i} = \int_\Omega w(x) \left(
-g_2(x)\,\frac{\partial N_i}{\partial x_1}
+ g_1(x)\,\frac{\partial N_i}{\partial x_2}
\right) \dd V .
\label{eq:load_psi}
\end{equation}
The entire infrastructure for assembling and solving the potential system is
reused verbatim; the only additional cost is the assembly of $F_\psi$ and a
second forward--backward substitution with the already-factorised stiffness
matrix.

\subsection{Recovering the non-conservative force}

From the relation $g = -D^{-1} f$ and $f = -D\,\grad U + v$, we obtain
$g = \grad U - D^{-1} v$, so that the solenoidal residual of the effective
field is $r = g - \grad U = -D^{-1} v$. The non-conservative force is therefore
recovered as
\begin{equation}
v = -D \cdot r = -D \cdot (g - \grad U) = f + D\,\grad U .
\label{eq:v_recovery}
\end{equation}
This relation is exact and does not require solving for $\psi$. The stream
function provides a smooth, basis-expanded representation of the same
solenoidal component, useful for visualisation and for computing the
reconstructed effective field
\begin{equation}
g_{\mathrm{recon}} = \grad U + \perpgrad \psi,
\end{equation}
which satisfies $\norm{g - g_{\mathrm{recon}}}_{L^2} = 0$ up to discretisation
error.

\subsection{Extension to higher dimensions}

In $D = 3$ the scalar stream function is replaced by a vector potential
$A \in H(\mathrm{curl};\Omega)$, yielding $g = \grad U + \curl A + h$, where
$h$ is a harmonic field. On simply connected domains $h = 0$; on multiply
connected domains the dimension of the harmonic space equals the first Betti
number $\beta_1$ \cite{huang2016}. A gauge condition (typically the Coulomb
gauge $\Div A = 0$) must be imposed, requiring a Lagrange-multiplier or
augmented-Lagrangian formulation \cite{wang2016}. For $D \ge 4$ the proper
framework is the Hodge decomposition of differential forms
\cite{hirani2003,desbrun2005}; the coexact component derives from a 2-form with
$\binom{n}{2}$ independent components, and tensor-product B-spline IgA becomes
intractable due to the exponential scaling of degrees of freedom. The dominant
practical approach, exemplified by the ddHodge framework \cite{ddhodge2025}, is
to apply dimensionality reduction before performing the Hodge decomposition.

%==============================================================================
\section{Estimation of Kramers--Moyal coefficients with P-splines}
\label{sec:psplines}
%==============================================================================

\subsection{From local kernel regression to global penalised regression}

In the original pipeline the conditional moments \eqref{eq:km_moments} are
estimated pointwise by Nadaraya--Watson kernel regression
\cite{wand1995,rydin2019} on a regular grid of bin centres, using the
\texttt{kramersmoyal} library. That estimator is \emph{local}: each evaluation
point aggregates the increments falling inside a kernel window of bandwidth
$h$. This design has three structural costs: (i) the data must be implicitly
binned and the estimates are only available at bin centres; (ii) near the edges
of the data domain the kernel leaks outside the support, producing a systematic
boundary bias that must be corrected by post-smoothing; (iii) the estimator has
several coupled free parameters---kernel shape, bandwidth, bin count,
post-smoothing scale, and density threshold---whose calibration is nontrivial.

Here we adopt the alternative suggested for this project: estimate both fields
\emph{globally} by P-splines \cite{eilers1996}. Each component of the drift and
each entry of the diffusion matrix is modelled as a tensor-product B-spline
surface whose coefficients are fitted by penalised least squares, with a
discrete difference penalty that controls smoothness. No spatial
discretisation of the data is involved: every observed transition
$(X_i \to X_{i+1})$ contributes directly to the regression at its exact
location $X_i$.

\subsection{Tensor-product B-spline basis}

Let the estimation domain be the box
$\Omega_{\mathrm{est}} = \prod_{d=1}^{D} [a_d, b_d]$, determined from the data
by coordinate-wise quantiles with padding (Section~\ref{sec:domain}). On each
axis we place an open, uniform knot vector with $n_{\text{basis}}$ B-spline
basis functions $N^{(d)}_{j,p}$ of degree $p$ (Cox--de Boor recursion
\cite{piegl2012}). The tensor-product basis is
\begin{equation}
N_{\mathbf{j}}(x) = \prod_{d=1}^{D} N^{(d)}_{j_d, p}(x_d),
\qquad
\mathbf{j} \in \mathcal{J},\quad K = |\mathcal{J}| = \prod_{d=1}^{D} n_{\text{basis}} ,
\label{eq:tensor_basis}
\end{equation}
and every scalar field is expanded as $\phi_h(x) = \sum_{\mathbf{j}}
c_{\mathbf{j}} N_{\mathbf{j}}(x)$. The total number of coefficients $K$ grows
exponentially with $D$; a guard (\texttt{max\_tensor\_bases}) aborts the fit
when $K$ exceeds a configurable budget.

\subsection{Penalised least squares and the difference penalty}

Smoothness is imposed through the P-spline penalty of Eilers and Marx
\cite{eilers1996}: along each axis, the $m$-th order differences of the
coefficients are penalised. For axis $d$ let $\Delta_m^{(d)}$ denote the
$m$-th order difference matrix acting on the coefficient index $j_d$. The
penalty matrix is
\begin{equation}
P = \sum_{d=1}^{D} I \otimes \cdots \otimes
\underbrace{(\Delta_m\topT \Delta_m)}_{\text{axis } d} \otimes \cdots \otimes I ,
\label{eq:penalty}
\end{equation}
so that $\mathbf{c}\topT P\, \mathbf{c}$ measures the roughness of the fitted
surface. With penalty order $m = 2$ (the default) the limit $\lambda \to
\infty$ drives the fit toward a global polynomial of degree $m-1$ (a plane),
while $\lambda \to 0$ approaches an interpolant of the data.

\subsection{Drift estimation by regression on the increments}

For the first-order coefficient we use the finite-$\Delta t$ identity
$\mathbb{E}[\Delta X \mid X_t = x] \approx f(x)\,\Delta t$ and regress the
scaled increments directly: for each component $d$ we solve
\begin{equation}
\min_{\mathbf{c}^{(d)}}
\sum_{i=1}^{N-1}
\left(
\frac{(X_{i+1} - X_i)_d}{\Delta t}
- \sum_{\mathbf{j}} c^{(d)}_{\mathbf{j}}\, N_{\mathbf{j}}(X_i)
\right)^{2}
+ \lambda_{\mathrm{drift}}\, (\mathbf{c}^{(d)})\topT P\, \mathbf{c}^{(d)} .
\label{eq:drift_fit}
\end{equation}
This is a ridge-type linear system
$(B\topT B + \lambda_{\mathrm{drift}} P)\, \mathbf{c}^{(d)} = B\topT y^{(d)}$,
where $B$ is the (sparse, row-wise) basis evaluation matrix and
$y^{(d)}_i = (X_{i+1} - X_i)_d / \Delta t$. All $D$ systems share the same
normal matrix, so the factorisation is computed once and reused across
components. Because the regression uses every transition at its exact
location, there is no binning error and no boundary truncation: the penalty
controls the extrapolation behaviour where data are scarce.

\subsection{Diffusion estimation and enforcement of positive-definiteness}
\label{sec:diffusion_fit}

For the second-order coefficient we first form the residuals of the drift fit,
$R_i = \Delta X_i - \Delta t\, \hat f(X_i)$, and regress the matrix-valued
targets
\begin{equation}
S_i = \frac{R_i R_i\topT}{2\,\Delta t}
\label{eq:diff_target}
\end{equation}
entry by entry: for each $(d, e)$ with $d \le e$ we solve a penalised system
analogous to \eqref{eq:drift_fit} with penalty weight
$\lambda_{\mathrm{diffusion}}$, and symmetrise by
$\hat D_{ed} = \hat D_{de}$. This \emph{method of moments} route guarantees
consistency between the estimated drift and diffusion.

The pointwise estimate $\hat D(x)$ is symmetric by construction but not
guaranteed positive semi-definite under penalisation. We enforce SPD in two
stages: (i) a floor on the diagonal entries (\texttt{min\_diffusion}) and a
cap on the implied correlations (\texttt{correlation\_clip}), applied to the
fitted coefficients; (ii) a final spectral check at evaluation time via
\texttt{np.linalg.eigh}, clipping negative eigenvalues. The diagnostics report
the minimum eigenvalue and condition number of $\hat D$ over a probe grid, so
that violations are visible rather than silently corrected.

\subsection{Estimation domain and output grid}
\label{sec:domain}

The estimation domain $\Omega_{\mathrm{est}}$ is defined coordinate-wise from
the data: with quantile level $q$ (\texttt{domain\_quantile}, default $0.005$)
and relative padding $\delta$ (\texttt{domain\_padding}, default $0.05$),
\begin{equation}
a_d = Q_d(q) - \delta\, (Q_d(1-q) - Q_d(q)),
\qquad
b_d = Q_d(1-q) + \delta\, (Q_d(1-q) - Q_d(q)),
\label{eq:domain}
\end{equation}
where $Q_d(\cdot)$ is the empirical quantile function of coordinate $d$.
Quantile trimming makes the domain robust to outlier excursions that would
otherwise stretch the knot vectors and waste basis functions.

The estimator delivers the same interface as the kernel pipeline: given a
requested output resolution \texttt{bins} $= (n_1, \dots, n_D)$ it returns the
fields evaluated on a regular grid of cell centres, plus the grid edges. By
default the output grid is \emph{inset} by a relative margin
(\texttt{grid\_margin}, default $0.30$): the outer $30\%$ of the estimation
domain on each side is excluded from the output grid. The rationale is
downstream: the IgA stage builds a density mask on the output grid, and if the
mask touches the grid boundary the Helmholtz solver develops spurious boundary
modes (with amplitudes orders of magnitude above the physical range). The
margin guarantees an empty frame between the data support and the grid edge,
so the mask never touches the boundary. This parameter is part of the
estimation configuration because it determines where the P-spline surfaces are
evaluated.

\subsection{The penalisation parameter as the analogue of the bandwidth}

The penalisation weights $\lambda_{\mathrm{drift}}$ and
$\lambda_{\mathrm{diffusion}}$ control the bias--variance trade-off exactly as
the bandwidth $h$ does in the kernel estimator: small $\lambda$ follows the
data closely (low bias, high variance); large $\lambda$ oversmooths (high
bias, low variance). Two structural differences are worth noting. First,
$\lambda$ acts \emph{globally} and dimensionlessly on the coefficient
roughness, so a single value adapts to varying local data density---dense
regions are fitted accurately while sparse regions are automatically smoothed
toward the penalty limit, without the occupancy failures of fixed-bandwidth
kernels. Second, the resolution of the estimator is set independently by
$n_{\text{basis}}$, decoupling the two choices that the bin count entangles in
the kernel pipeline. Section~\ref{sec:hyperparameters} defines the calibration
strategy for $\lambda$ and the remaining hyperparameters.

\subsection{Implementation and interface}

The estimator is implemented in
\texttt{src/potential\_reconstruction/pspline\_estimator.py} and is integrated
into the pipeline through a dispatch parameter: \texttt{extract\_km\_coefficients}
accepts \texttt{estimation\_method='nadaraya\_watson'} (default, historical
behaviour, unchanged) or \texttt{estimation\_method='p\_splines'}, with the
hyperparameters passed as a configuration dictionary. In P-splines mode the
legacy parameters ($p$, \texttt{bw}, \texttt{kernel}, \texttt{sigma\_smooth},
\texttt{density\_threshold}) do not apply; masking by statistical support is
left to the IgA stage, which already weights cells by empirical density. The
call returns \texttt{(drift, diffusion, edges)} as before, and optionally a
diagnostics dictionary (fit RMSE and $R^2$ per drift component, minimum
eigenvalue and condition number of $\hat D$, basis and penalty configuration).

%==============================================================================
\section{Reconstruction via isogeometric analysis}
%==============================================================================

The reconstruction stage is identical to the original pipeline; we summarise it
here for completeness.

\subsection{Robust inversion of the diffusion matrix}

The pointwise estimate $\hat D(x_k)$ at each cell may be anisotropic or nearly
singular. We employ the Moore--Penrose pseudo-inverse via SVD with adaptive
truncation \cite{golub2013}. At each cell $k$,
$\hat D_k = V_k \Sigma_k V_k\topT$, and the truncation threshold
$\tau_k = \max(\texttt{atol},\, \texttt{rtol}\, \sigma_{k,1})$ (typically
$\texttt{rtol} = 10^{-6}$, $\texttt{atol} = 10^{-12}$) nulls directions with
insufficient statistical information, so that
$g_k = -\hat D_k^{+} f_k$ is determined along those axes by the global
smoothness of the variational solver rather than by amplifying noise.

\subsection{Tensor-product B-spline space}

Both the potential and the stream function are expanded in tensor-product
B-splines \cite{piegl2012,cottrell2009}:
\begin{equation}
U_h(x) = \sum_{j \in J} c_j N_j(x),
\qquad
\psi_h(x) = \sum_{j \in J} d_j N_j(x),
\end{equation}
with partition of unity, compact support, non-negativity, and $C^{p-1}$
continuity at simple knots. For $p = 3$ the potential is $C^2$ and its
gradient $C^1$, eliminating the need for arbitrary post-hoc smoothing.

\subsection{Galerkin formulation and system assembly}

Substituting $U_h$ into the weak form \eqref{eq:variational} and taking test
functions $v = N_i$ (Bubnov--Galerkin \cite{cottrell2009}) yields the linear
system $K\,\mathbf{c} = F_U$ with
\begin{equation}
K_{ij} = \sum_{d=1}^{D} \int_\Omega w(x)\,
\frac{\partial N_i}{\partial x_d} \frac{\partial N_j}{\partial x_d} \, \dd V,
\qquad
F_{U,i} = \sum_{d=1}^{D} \int_\Omega w(x)\, g_d(x)\,
\frac{\partial N_i}{\partial x_d} \, \dd V,
\end{equation}
where the weight function
$w(x) = \min\big(1,\, \rho(x)/\rho_{\mathrm{threshold}}\big)$,
$\rho_{\mathrm{threshold}} = 0.05 \max_x \rho(x)$, weights each cell by its
statistical quality. For the stream function the same stiffness matrix is
reused with the load vector \eqref{eq:load_psi}.

\begin{theorem}
Let $U$ be the solution of $-\grad^2 U = -\Div g$ with natural Neumann
conditions $\partial U / \partial n = g \cdot n$. Then $\grad U$ is the
orthogonal projection of $g$ onto the gradient subspace $\mathcal{G}$ in the
$L^2$ sense.
\end{theorem}
\begin{proof}
The weak form reads $\int_\Omega \grad U \cdot \grad v \, \dd V = \int_\Omega g
\cdot \grad v \, \dd V$ for all $v \in H^1(\Omega)$, i.e.\
$\int_\Omega (\grad U - g) \cdot \grad v \, \dd V = 0$, which is precisely the
orthogonality condition $(\grad U - g) \perp \mathcal{G}$. Since $\grad U \in
\mathcal{G}$ by construction, $r = g - \grad U \in \mathcal{G}^{\perp} =
\mathcal{S}$.
\end{proof}

\subsection{Solver and potential anchoring}

Each system has rank $N_{\mathrm{dof}} - 1$, with kernel the constant
functions; we fix the coefficient whose basis function peaks at the cell of
highest empirical density via the penalty method (a penalty of $10^{12}
\max|K_{ij}|$ on the diagonal). For $N_{\mathrm{dof}} < 5000$ we use a direct
supernodal sparse LDL$\topT$ factorisation \cite{golub2013}; for larger
systems, preconditioned conjugate gradients with incomplete Cholesky
preconditioning. For the stream-function system the factorisation of $K$ is
reused.

%==============================================================================
\section{Pipeline synthesis}
%==============================================================================

The complete reconstruction workflow integrates the stages described above.
The input is a latent trajectory $x(t) \in \R^D$ extracted from EEG by PCA
whitening and ICA. First, the Kramers--Moyal stage estimates the drift
$\Drift$ and diffusion $\Diff$ fields---now by global P-spline regression
(Section~\ref{sec:psplines}), or optionally by the historical
Nadaraya--Watson kernel estimator, selected by a single configuration
parameter---together with the empirical density $\rho(x)$ and the grid edges.
Second, the effective field $g = -D^{-1} f$ is computed via robust SVD
inversion. Third, the IgA stage reconstructs the scalar potential $U(x)$ and,
in two dimensions, the stream function $\psi(x)$ via the decoupled Poisson
system, from which the solenoidal residual, the non-conservative force
$v = -D \cdot r$, the reconstructed effective field $g_{\mathrm{recon}}$, and
the non-equilibrium metric $\eta$ are derived. The pipeline is fully
observable: the P-spline diagnostics (fit RMSE, $R^2$, minimum eigenvalue and
condition number of $\hat D$) and the local IgA diagnostics
($\mathrm{rank}_D(x)$, $\mathrm{cond}_D(x)$) flag unreliable regions, while
$\eta$ guides the physical interpretation without requiring any ground truth.

% [PLACEHOLDER: Figura del diagrama de flujo del pipeline de dos etapas
%  (análoga a la Fig.~3 del reporte original), actualizada para indicar el
%  despacho NW / P-splines en la etapa de estimación.]

%==============================================================================
\section{Hyperparameters and calibration strategy}
\label{sec:hyperparameters}
%==============================================================================

The kernel pipeline required calibrating the bandwidth $h$ (via the hybrid
score described in the original report), plus the kernel shape, bin count,
post-smoothing scale and density threshold. The P-spline estimator replaces
this set with the following hyperparameters, defined in the
\texttt{PSplineConfig} configuration object (defaults in parentheses):

\begin{table}[h]
\centering
\small
\begin{tabular}{@{}llp{6.8cm}@{}}
\toprule
Parameter & Default & Role \\
\midrule
\texttt{n\_basis} & 8 & B-spline basis functions per axis; sets the spatial
resolution of the estimator (total coefficients $K = n_{\text{basis}}^D$). \\
\texttt{degree} & 3 & Degree $p$ of the B-spline basis ($C^{p-1}$ fields). \\
\texttt{penalty\_order} & 2 & Order $m$ of the difference penalty; the
large-$\lambda$ limit is a polynomial of degree $m-1$. \\
\texttt{lambda\_drift} & $10^{-3}$ & Penalisation weight of the drift fit;
the analogue of the kernel bandwidth for $\Drift$. \\
\texttt{lambda\_diffusion} & $10^{-3}$ & Penalisation weight of the diffusion
fit; the analogue of the bandwidth for $\Diff$. \\
\texttt{domain\_quantile} & $0.005$ & Quantile level defining the estimation
domain (outlier trimming). \\
\texttt{domain\_padding} & $0.05$ & Relative padding of the estimation domain
beyond the quantile range. \\
\texttt{grid\_margin} & $0.30$ & Relative inset of the output grid; keeps the
IgA density mask away from the boundary (prevents spurious modes). \\
\texttt{min\_diffusion} & $10^{-12}$ & Floor on the diagonal of $\hat D$
(SPD enforcement). \\
\texttt{correlation\_clip} & $0.999$ & Cap on implied correlations of
$\hat D$ (SPD enforcement). \\
\bottomrule
\end{tabular}
\caption{Hyperparameters of the P-spline Kramers--Moyal estimator.}
\label{tab:hyperparameters}
\end{table}

The parameters that genuinely require \emph{optimisation}---in the same sense
that the bandwidth required it in the kernel pipeline---are the penalisation
weights \texttt{lambda\_drift} and \texttt{lambda\_diffusion}, and secondarily
\texttt{n\_basis}. The remaining parameters are structural: \texttt{degree}
and \texttt{penalty\_order} fix standard smoothness classes; the domain and
margin parameters are set by robustness considerations rather than by fit
quality; the SPD parameters are numerical guards.

We propose the following calibration strategy, analogous in spirit to the
bandwidth scan of the kernel pipeline but cheaper, because each candidate fit
reuses the same basis matrix:

\begin{enumerate}
\item \textbf{Candidate grid.} Scan $\lambda$ on a log-spaced grid spanning
interpolation to near-parametric behaviour (e.g.\ $10^{-6}$ to $10^{2}$), at
fixed \texttt{n\_basis}.
\item \textbf{Score each candidate} on the reconstructed potential with the
same ground-truth-free hybrid score used for the bandwidth: (a) correlation
of $U$ with $-\log\rho_{\mathrm{emp}}$ over the density mask; (b) stability of
$U$ between neighbouring $\lambda$ values; (c) plausibility of the range of
$U$ relative to the range of $-\log\rho$.
\item \textbf{Statistical cross-checks} from the estimator diagnostics: drift
$R^2$ and RMSE per component, and the minimum eigenvalue and condition number
of $\hat D$ on a probe grid. Candidates producing near-singular or negative
diffusion spectra are rejected regardless of score.
\item \textbf{Resolution check.} Repeat at two or three values of
\texttt{n\_basis}; retain the smallest basis whose best-$\lambda$ score
plateaus, to avoid the exponential cost $K = n_{\text{basis}}^D$ in high
dimension.
\end{enumerate}

% [PLACEHOLDER: figura del escaneo de lambda (score híbrido vs lambda, con el
%  óptimo marcado), análoga al panel de optimal_bw del pipeline NW. Pendiente
%  de la implementación del scan en bw_optimization para P-splines.]

%==============================================================================
\section{Numerical experiments}
\label{sec:experiments}
%==============================================================================

% [PLACEHOLDER: sección de resultados del método P-splines. Estructura prevista,
%  idéntica a la del reporte original para permitir comparación directa:
%
%  7.1 Ornstein–Uhlenbeck (2D)      — sistema de equilibrio, potencial cuadrático
%  7.2 Double-well potential (2D)   — equilibrio, dos pozos con acoplamiento
%  7.3 Multi-stable potential (2D)  — equilibrio, cuatro pozos
%  7.4 Ring attractor (2D)          — NO equilibrio, componente rotacional, eta >> 0.1
%  7.5 Triple-well potential (3D)   — caso de mayor dimensión
%
%  Para cada sistema: figura de campos drift/diffusion estimados vs teóricos,
%  figura del potencial reconstruido vs teórico, correlación de forma (Pearson)
%  y RMSE normalizado, métrica eta, y diagnósticos P-splines (lambda usado,
%  R2 del drift, eig/cond de D). En 7.4 además: fuerza no-conservativa v,
%  stream function psi y visualización combinada U + v + streamlines.
%
%  Los resultados se generarán con src/pipelines/test_iga_from_models.py
%  fijando estimation_method='p_splines' en la configuración de cada modelo.
%  NOTA: no incluir cifras de resultados fuera de esta sección.]

%==============================================================================
\section{Open problems and future directions}
%==============================================================================

Several directions remain open and are the subject of ongoing work.

\textbf{Automatic penalisation selection.} The scan-based calibration of
Section~\ref{sec:hyperparameters} mirrors the bandwidth selection of the
kernel pipeline. A more principled alternative is generalised cross-validation
or REML estimation of $\lambda$ within the penalised regression itself, which
would make the estimator fully self-calibrating.

\textbf{3D rotational component.} The current implementation computes the
stream function $\psi$ only in two dimensions. Extending the pipeline to three
dimensions requires solving for a vector potential $A$ subject to the Coulomb
gauge $\Div A = 0$ \cite{huang2016,wang2016}. For $D \ge 4$ the
dimensionality-reduction paradigm offers a computationally viable alternative
\cite{ddhodge2025}.

\textbf{Curse of dimensionality in the estimator.} The tensor-product basis
scales as $K = n_{\text{basis}}^D$ coefficients. Sparse-grid or
low-rank-tensor variants of the P-spline fit would extend the estimator to
$D \ge 4$ without the current \texttt{max\_tensor\_bases} guard aborting the
fit.

\textbf{State-dependent diffusion.} The SVD-based inversion handles
anisotropic $D(x)$ cell by cell, but the variational formulation treats $D$ as
given. Jointly estimating $U(x)$ and $D(x)$ within the same Galerkin framework
would unify the reconstruction into a single self-consistent problem.

%==============================================================================
\begin{thebibliography}{99}

\bibitem{risken1989}
Risken, H. (1989). \emph{The Fokker-Planck Equation: Methods of Solution and
Applications}. Springer.

\bibitem{tabar2019}
Tabar, M.~R.~R. (2019). \emph{Analysis and Data-Based Reconstruction of
Complex Nonlinear Dynamical Systems}. Springer.

\bibitem{gardiner2009}
Gardiner, C.~W. (2009). \emph{Stochastic Methods: A Handbook for the Natural
and Social Sciences}, 4th ed. Springer.

\bibitem{silverman1986}
Silverman, B.~W. (1986). \emph{Density Estimation for Statistics and Data
Analysis}. Chapman \& Hall.

\bibitem{wand1995}
Wand, M.~P., Jones, M.~C. (1995). \emph{Kernel Smoothing}. Chapman \& Hall.

\bibitem{epanechnikov1969}
Epanechnikov, V.~A. (1969). Non-parametric estimation of a multivariate
probability density. \emph{Theory of Probability and Its Applications}, 14(1),
153--158.

\bibitem{golub2013}
Golub, G.~H., Van Loan, C.~F. (2013). \emph{Matrix Computations}, 4th ed.
Johns Hopkins University Press.

\bibitem{strang1986}
Strang, G. (1986). \emph{Introduction to Applied Mathematics}.
Wellesley-Cambridge Press.

\bibitem{rydin2019}
Rydin Gorj\~ao, L. et al. (2019). kramersmoyal: Kramers--Moyal coefficients
for stochastic processes. \emph{Journal of Open Source Software}.

\bibitem{piegl2012}
Piegl, L., Tiller, W. (2012). \emph{The NURBS Book}, 2nd ed. Springer.

\bibitem{cottrell2009}
Cottrell, J.~A., Hughes, T.~J.~R., Bazilevs, Y. (2009). \emph{Isogeometric
Analysis: Toward Integration of CAD and FEA}. Wiley.

\bibitem{helmholtz1858}
Helmholtz, H. von (1858). \"Uber Integrale der hydrodynamischen Gleichungen,
welche den Wirbelbewegungen entsprechen. \emph{Journal f\"ur die reine und
angewandte Mathematik}, 55, 25--55.

\bibitem{auchmuty2016}
Auchmuty, G. (2016). Bounds and regularity of solutions of planar div-curl
problems. \emph{arXiv preprint arXiv:1610.06827}.

\bibitem{huang2016}
Huang, X. (2016). Differential complexes, Helmholtz decompositions, and
decoupling of mixed methods. \emph{arXiv preprint arXiv:1611.03936}.

\bibitem{wang2016}
Wang, J., Ye, X. (2016). Discretization of div-curl systems by weak Galerkin
finite element methods. \emph{Journal of Scientific Computing}, 68,
1144--1161.

\bibitem{hirani2003}
Hirani, A.~N. (2003). \emph{Discrete Exterior Calculus}. PhD thesis,
California Institute of Technology.

\bibitem{desbrun2005}
Desbrun, M. et al. (2005). Discrete exterior calculus for variational
problems. \emph{arXiv preprint math/0508341}.

\bibitem{ddhodge2025}
ddHodge Framework Authors (2025). ddHodge: A framework for approximating the
Hodge components on a low-dimensional manifold to study high-dimensional
dynamics of cell states. \emph{bioRxiv}. doi:10.1101/2025.02.20.639398.

\bibitem{eilers1996}
Eilers, P.~H.~C., Marx, B.~D. (1996). Flexible smoothing with B-splines and
penalties. \emph{Statistical Science}, 11(2), 89--121.

\end{thebibliography}

\end{document}
